\documentclass[10pt,aspectratio=169]{beamer}

\usetheme{metropolis}
\usepackage[utf8]{inputenc}
\usepackage[T1]{fontenc}
\usepackage[english]{babel}
\usepackage{listings}
\usepackage{booktabs}
\usepackage{tabularx}
\usepackage{amsmath,amssymb}

\definecolor{gfblue}{HTML}{1F4E79}
\definecolor{codebg}{HTML}{F4F4F2}
\definecolor{warn}{HTML}{B03A2E}
\definecolor{good}{HTML}{1E6B3A}

\setbeamercolor{frametitle}{bg=gfblue}
\setbeamercolor{alerted text}{fg=warn}

\lstset{
  basicstyle=\ttfamily\scriptsize,
  backgroundcolor=\color{codebg},
  frame=none,
  columns=fullflexible,
  keepspaces=true,
  showstringspaces=false,
  breaklines=true,
  xleftmargin=2mm,
  literate=
    {é}{{\'e}}1 {è}{{\`e}}1 {ê}{{\^e}}1 {à}{{\`a}}1
    {→}{{$\rightarrow$}}1 {∀}{{$\forall$}}1 {∃}{{$\exists$}}1 {∈}{{$\in$}}1
    {ℕ}{{$\mathbb{N}$}}1 {∣}{{$\mid$}}1 {∨}{{$\vee$}}1 {∧}{{$\wedge$}}1
    {≠}{{$\neq$}}1 {≤}{{$\leq$}}1 {⟨}{{$\langle$}}1 {⟩}{{$\rangle$}}1
    {↔}{{$\leftrightarrow$}}1 {λ}{{$\lambda$}}1
}

\newcommand{\bad}[1]{\textcolor{warn}{#1}}
\newcommand{\ok}[1]{\textcolor{good}{#1}}
\newcommand{\commit}[1]{\texttt{\footnotesize #1}}
\newcommand{\prompt}[1]{\begin{quote}\small\itshape ``#1''\end{quote}}
\newcommand{\decided}[1]{\par\textcolor{gfblue}{\textbf{#1}}\par}

\title{Informalizing Lean Mathlib}
\subtitle{Mathlib in English, French, Czech, German and Swedish with Informath, built with Claude Code}
\author{Aarne Ranta}
\date{23 October 2026}

\begin{document}

\maketitle

% ------------------------------------------------------------------ background

\begin{frame}[fragile]{Lean and Mathlib}
  \small
  \begin{columns}[T]
    \begin{column}{0.5\textwidth}
      \begin{itemize}
        \item \textbf{Lean 4}: a proof assistant and functional programming
              language based on dependent type theory (de Moura et al.;
              now the Lean FRO).
        \item \textbf{Mathlib}: Lean's community library of formalized
              mathematics, from algebra and number theory to analysis,
              topology and category theory.
        \item At the commit used here (\commit{c31114ff81}, 21 Oct 2026):
              8,650 files, 2.3 million lines, about 194,000 theorems and
              48,000 definitions, structures and classes.
        \item Statements are read by Lean's \emph{elaborator}: type classes,
              implicit arguments, coercions and notation are resolved there,
              not in the source text.
      \end{itemize}
    \end{column}
    \begin{column}{0.48\textwidth}
\begin{lstlisting}
/-- `Nat.Prime p` means that `p` is a prime number,
  that is, a natural number at least 2 whose only
  divisors are `p` and `1`. -/
@[pp_nodot, wikidata Q49008, oeis A000040]
def Prime (p : ℕ) :=
  Irreducible p

theorem exists_prime_and_dvd {n : ℕ} (hn : n ≠ 1) :
    ∃ p, Prime p ∧ p ∣ n :=
  ⟨minFac n, minFac_prime hn, minFac_dvd _⟩
\end{lstlisting}
      \footnotesize\centering\texttt{Mathlib/Data/Nat/Prime/Defs.lean}
    \end{column}
  \end{columns}
\end{frame}

\begin{frame}{The problem: Mathlib is formal}
  \small
  \textbf{What informal text there is}
  \begin{itemize}
    \item \textbf{Docstrings}, written by hand. Batteries' \texttt{docBlame}
          linter requires them for definitions; \texttt{docBlameThm}, the
          same for theorems, is disabled. In our checkout,
          \textbf{83\%} of definitions, structures and classes have one, but only
          \textbf{15\%} of theorems and lemmas.
    \item \textbf{Module docstrings} summarizing each file, and the HTML
          documentation generated by doc-gen4, which shows the formal
          statements with the docstrings.
    \item \textbf{LLM informalizations} made for search and for training:
          LeanSearch (Gao et al., EMNLP Findings 2024) and the Herald
          dataset (Gao et al., ICLR 2025), checked by back-translation.
  \end{itemize}
  \textbf{What is missing}
  \begin{itemize}
    \item Readings of \emph{all} statements that are faithful by construction,
          reproducible, and in several languages.
    \item Proofs in natural language: none of the above covers them.
  \end{itemize}
\end{frame}

\begin{frame}{GF and Informath}
  \small
  \begin{columns}[T]
    \begin{column}{0.5\textwidth}
      \textbf{GF}, Grammatical Framework
      \begin{itemize}
        \item An abstract syntax (trees) with concrete syntaxes
              (linearizations) for many languages.
        \item The Resource Grammar Library takes care of morphology and
              syntax: agreement, word order, inflection.
      \end{itemize}
      \textbf{Informath} (2024--, 932 commits)
      \begin{itemize}
        \item Translates between \textbf{Dedukti} and natural language,
              in both directions, with LaTeX formulas.
        \item A Dedukti dialect: \texttt{Set}, \texttt{Elem}, \texttt{Prop},
              \texttt{Proof}, and a base of arithmetic and logic.
        \item \textbf{Symbol tables} (\texttt{.dkgf}) give each constant
              its verbal and symbolic forms.
      \end{itemize}
    \end{column}
    \begin{column}{0.48\textwidth}
      \textbf{From Dedukti to text}
      \begin{enumerate}
        \item Annotate the constants with the symbol tables.
        \item Translate to \emph{MathCore}, a minimal GF syntax.
        \item Generate the variants: verbal or symbolic, aggregated,
              in-situ quantifiers\ldots
        \item Linearize them in English, French, Swedish, \ldots
              or as LaTeX formulas.
        \item Rank them, and pick the best.
      \end{enumerate}
      \textbf{Proof texts} (\texttt{-proof-text}): numbered lines, each
      justified by a rule or a lemma, as in godement-index.
    \end{column}
  \end{columns}
\end{frame}

% ------------------------------------------------------------------ the process

\section{Building mathlib-informath}

\begin{frame}{Two days, thirteen prompts}
  \small
  \scriptsize\begin{tabular}{@{}r>{\raggedright\arraybackslash}p{0.42\textwidth}>{\raggedright\arraybackslash}p{0.3\textwidth}l@{}}
    \toprule
    & prompt (abridged) & outcome & commits \\
    \midrule
    1 & build an informalization system for Mathlib & exporter, Nat primes & \commit{c46c6c3} \\
    2 & build mathlib.dkgf and apply it & 17 entries: 69 $\to$ 7 unworded & \\
    3 & an HTML page with Lean, Dedukti, MathJax, GitHub links & the page, then GitHub Pages & \\
    4 & split it into subfiles; Factorization next & 322 statements on 2 pages & \commit{3e8bc29} \\
    5 & keep the definientia, show the proofs & 84 proof texts; Informath fix & \commit{5e78612} \commit{1658da6} \\
    6 & natural language from the sources & docstrings and names; \texttt{Mathlib.gf} & \commit{8244cae} \\
    7 & macros in red: MathJax issue? & a builder bug; $\operatorname{ord}_p(n)$ & \commit{6b7383c} \\
    8 & the Algebra and Order.Defs modules & 705 statements on 4 pages & \commit{976bbd0} \\
    9 & add French & a third row, an index & \commit{823155a} \\
    10 & add Czech, hidden with French; buttons at each entry & a fourth row & \commit{cacc011} \commit{0c9c416} \\
    11 & aggregate hypotheses (in Informath); rebuild the pages & 503 readings change & \commit{afd0d9d} \\
    12 & add German & a fifth row, with a time limit & \commit{a53c372} \\
    13 & proofs that use tactics, as proof units & 25 tactic proofs (English) & \commit{9822ffa} \\
    \bottomrule
  \end{tabular}

  \vspace{2mm}
  \footnotesize 21--22 October 2026: 13 commits in mathlib-informath (besides these slides)
  and 8 in informath (\commit{4f31f55}, \commit{ceb9b85}, \commit{7773194},
  \commit{778e998}, \commit{0da10d2}, \commit{5907eec}, \commit{a8b863f}, \commit{4bcff16}),
  all made by Claude when asked, and pushed by me.
\end{frame}

\begin{frame}{1. The plan: three decisions}
  \prompt{you can see ../mathlib4, right? It is Leans famous Mathlib. I would
  like to build an informalization system for it using Informath.}
  Claude surveyed Mathlib, Informath's Lean support and lean2dk (Deducteam's
  faithful Lean-to-Dedukti translator), then asked three questions:
  \begin{enumerate}
    \item \textbf{Input:} \ok{own Lean exporter} (recommended), lean2dk, or
          parsing the Lean source. lean2dk encodes the whole type theory,
          which is the opposite of what Informath reads; the source text
          lacks what the elaborator adds.
    \item \textbf{First slice:} Algebra/Group/Basic, \ok{Nat primes}, or
          orders and lattices.
    \item \textbf{Setup:} \ok{a new project, with Mathlib's build cache
          downloaded} (7 GB in two minutes, instead of hours of compiling).
  \end{enumerate}
  The answers are in green. The first run read all 130 statements of the five
  prime modules back into English, 61 of them already fully worded.
\end{frame}

\begin{frame}[fragile]{The exporter: from Lean's environment to Dedukti}
  \begin{columns}[T]
    \begin{column}{0.47\textwidth}
      \texttt{Export.lean} (now 500 lines) runs over Mathlib's elaborated
      environment:
      \begin{itemize}
        \item binders become \texttt{Elem}, \texttt{Proof} or \texttt{Set}
              hypotheses;
        \item implicit and instance arguments are dropped from applications;
        \item casts and coercions are the identity, numerals Informath's digits;
        \item \texttt{mapping.tsv} maps Lean's notions to Informath's base:
              \texttt{Nat.Prime}$\to$\texttt{prime}, \texttt{Dvd.dvd}$\to$\texttt{divide};
        \item per module: \texttt{.dk}, \texttt{.json} (docstrings, lines),
              \texttt{-deps.dk}.
      \end{itemize}
    \end{column}
    \begin{column}{0.51\textwidth}
\begin{lstlisting}
theorem Prime.eq_one_or_self_of_dvd {p : ℕ} (pp : p.Prime)
    (m : ℕ) (hm : m ∣ p) : m = 1 ∨ m = p
\end{lstlisting}
\begin{lstlisting}
Nat__Prime__eq_one_or_self_of_dvd : (p : Elem Nat) ->
  (pp : Proof (prime p)) -> (m : Elem Nat) ->
  (hm : Proof (divide m p)) -> Proof (or (Eq m 1) (Eq m p)).
\end{lstlisting}
      \small
      Let $p$ be a natural number. Assume that $p$ is prime. Let $m$ be a
      natural number. Assume that $m$ divides $p$. Then $m = 1$ or $m = p$.

      \vspace{1mm}
      $p \in \mathbb{N}, \mathrm{prime}(p), m \in \mathbb{N}, m \mid p
      \vdash m = 1 \vee m = p$
    \end{column}
  \end{columns}
\end{frame}

\begin{frame}{2. A symbol table for Mathlib's constants}
  \small
  \prompt{build mathlib.dkgf and apply it}
  \begin{itemize}
    \item 17 entries, worded after Mathlib's docstrings where the grammar had
          the words, checked against the exported types:
          \texttt{Nat.minFac} ``the minimal prime factor of \#1'',
          \texttt{Nat.Coprime} ``\#1 and \#2 are coprime'' $|$ $\#1 \perp \#2$.
    \item Readings with an unworded constant: \textbf{69 $\to$ 7} of 130.
  \end{itemize}
  \decided{Writing the table exposed three exporter bugs}
  \begin{itemize}
    \item Prop-valued constants were typed \texttt{Set}.
    \item Lean's \texttt{Set} clashed with Informath's: it now maps to \texttt{powerset}.
    \item Definitions were declared with implicit arguments that the
          applications drop, so \texttt{\#k} counted the wrong arguments.
  \end{itemize}
  \decided{And Informath issues, reported but not fixed}
  \texttt{\$} inside a quoted verbal form splits a table entry; a proposition
  inside a macro argument loses its notation; ``does not have \bad{a} upper bound''.
\end{frame}

\begin{frame}{3. The page, and GitHub Pages}
  \small
  \prompt{produce a mathlib-informath.html file, which shows for each statement
  the original lean, the dedukti translation, and the informalization in mathjax
  latex, as well as a link to the original lean source in mathlib github}
  \begin{itemize}
    \item Each entry: name, kind, docstring, the Lean source up to its
          \texttt{:=}, the Dedukti, English, the formula, and a link to the
          line at Mathlib's exact commit.
    \item The exporter now records source positions and keeps the source
          order. Lean-generated declarations, which have no position, are
          left out.
    \item Checked in headless Chrome: MathJax renders the readings.
  \end{itemize}
  \prompt{will it be hosted in github?}
  Not yet: the project had no remote. Claude explained the two steps,
  a repository and GitHub Pages, but left both to me, as they publish under
  my name. I created \texttt{aarneranta/mathlib-informath}; Claude committed
  \commit{c46c6c3}; I pushed.
\end{frame}

\begin{frame}{4. More of Mathlib: Factorization, and a page per part}
  \small
  \prompt{Do more of Mathlib. But now you should split the html document to a
  main file, which contains a table of contents with links to subfiles of
  reasonable size. \ldots\ Factorization is perhaps the next example?}
  \begin{itemize}
    \item One page per directory; the table of contents forwards old links
          (\texttt{\#Nat.Prime.two\_le}) to the right page, since the page is public.
    \item \textbf{Coercions to functions:} \texttt{n.factorization p} becomes one
          constant \texttt{Nat\_\_factorization\_\_app n p}.
    \item \textbf{Binders} ($\prod_{x \in s} f(x)$, $\{x \in s \mid P\}$): the
          exporter eta-expands the function arguments, and raw macros take the
          bound variable as a parameter, as Informath's \texttt{\textbackslash Summa}.
    \item \textbf{Type classes} become predicates: \texttt{[CommMonoid M]} reads
          ``Assume that $M$ is a commutative monoid''.
    \item About 30 new table entries; deprecated aliases left out; a
          \texttt{-force-symbolic} fallback for statements with no reading.
  \end{itemize}
  Result: 322 statements on two pages, 267 fully worded (\commit{3e8bc29}).
\end{frame}

\begin{frame}[fragile]{5. Definientia and proofs}
  \small
  \prompt{I don't think we need to cut judgements at := if the definiens is an
  explicit expression directly translatable to Dedukti. Keep it in those cases,
  and if the value type is Proof X, show the proof in the same way as in
  godement-index.}
  \begin{columns}[T]
    \begin{column}{0.5\textwidth}
      \begin{itemize}
        \item Only term definientia (not \texttt{by} tactics, not equations).
              Theorems' proofs needed \texttt{value? (allowOpaque := true)}.
        \item Proofs go into the natural deduction rules of \texttt{nd.dk}:
              $\lambda$s proving $\to$ or $\forall$ become \texttt{ifI}/\texttt{forallI},
              applied hypotheses \texttt{ifE}/\texttt{forallE}.
        \item Lemmas keep their arguments in proofs, except type class
              instances.
        \item An Informath fix (\commit{4f31f55}): \texttt{-proof-text} crashed on
              a $\lambda$ typed by a proof of an implication.
      \end{itemize}
    \end{column}
    \begin{column}{0.48\textwidth}
\begin{lstlisting}
theorem exists_prime_and_dvd {n : ℕ} (hn : n ≠ 1) :
    ∃ p, Prime p ∧ p ∣ n :=
  ⟨minFac n, minFac_prime hn, minFac_dvd _⟩
\end{lstlisting}
      \scriptsize
      \begin{tabular}{@{}r>{\raggedright\arraybackslash}p{0.6\linewidth}l@{}}
        1. & Let $n$ be an element of $\mathbb{N}$. & \\
        2. & Assume that $n \neq 1$. & \\
        3. & $\operatorname{minFac}(n)$ is prime. & \texttt{minFac\_prime}~2 \\
        4. & $\operatorname{minFac}(n) \mid n$. & \texttt{minFac\_dvd} \\
        5. & $\operatorname{minFac}(n)$ is prime and $\operatorname{minFac}(n) \mid n$. & \texttt{andI} 3, 4 \\
        6. & There exists an element $p$ of $\mathbb{N}$, such that $p$ is prime and $p \mid n$. & \texttt{existsI} 5 \\
      \end{tabular}
    \end{column}
  \end{columns}
\end{frame}

\begin{frame}{6. Natural language from the sources}
  \footnotesize
  \prompt{find more natural language expressions instead of verbatim Lean
  constants. \ldots\ Consider for instance primeFactorsList(n) \ldots\ Could
  this information be found somewhere in the available sources?}
  \begin{columns}[T]
    \begin{column}{0.5\textwidth}
      Three sources, as Claude found them:
      \begin{itemize}
        \item \textbf{docstrings}, in fixed patterns: ``\texttt{f n} is the X of
              \texttt{n}'', ``Returns the X\ldots'';
        \item \textbf{names}, by Mathlib's naming conventions:
              \texttt{IsUnit} $\to$ ``\#1 is a unit'';
        \item \textbf{Wikidata attributes}, but only on a few declarations.
      \end{itemize}
      \texttt{suggest-entries.py} makes candidates from the first two and keeps
      those that \texttt{-parse-example} parses.
    \end{column}
    \begin{column}{0.48\textwidth}
      \prompt{if you need new words, add them to Mathlib.gf}
      \begin{itemize}
        \item 31 reviewed entries: ``the smallest prime factor of \#1'',
              ``the set of prime factors of \#1''.
        \item \texttt{Mathlib.gf}: ``smallest'', ``descending'', ``bounded
              above'', ``set of divisors''\ldots\ (\commit{ceb9b85}).
        \item English is read from a table without pseudo-notations such as
              $\operatorname{primeFactors}(n)$, which the ranking preferred.
      \end{itemize}
      Factorization readings with a raw Lean name: \textbf{94 $\to$ 8} of 194.
    \end{column}
  \end{columns}
\end{frame}

\begin{frame}{7. Red macros}
  \prompt{In Factorization.html, some macros are rendered in red and not
  expanded. For instance, \textbackslash NatfactorizationappMACRoI. Is this a
  MathJax Issue or something easy to fix?}
  \begin{itemize}
    \item Not MathJax: Claude's own bug from step 6. Macros are numbered by
          their position in the table, and the English table, with fewer
          LaTeX forms, numbered them otherwise. 37 macros were undefined.
    \item The fix: the pages define the macros of both tables, with a check
          for clashing names (none).
    \item The screenshot check found another problem: the multiplicity of
          $p$ in $n$ and the $p$-adic valuation were written $\nu_p(n)$ and
          $v_p(n)$. Claude asked; I chose to commit both changes, and
          the multiplicity is now $\operatorname{ord}_p(n)$ (\commit{6b7383c}).
  \end{itemize}
\end{frame}

\begin{frame}{8. Down the imports: Algebra and Order}
  \small
  \prompt{I can see that already Prime.Defs imports some Algebra modules:
  Group, Prime, etc. Can you run informalization with those Algebra modules
  and also Order.Defs modules?}
  \begin{itemize}
    \item The five \texttt{Order.Defs} modules (354 statements) and two
          Algebra modules; \texttt{GroupWithZero.Nat} has only instances and
          was left out.
    \item \texttt{\# page: X} lines in \texttt{modules.txt} put small modules on
          one page.
    \item About 40 entries for order, algebra and relation classes: ``\#1 is
          partially ordered'', ``\#2 is a strict weak order'', ``the supremum of
          \#1 and \#2'' $|$ $\#1 \sqcup \#2$; four more words in \texttt{Mathlib.gf}.
    \item \texttt{Decidable} instance hypotheses are computational, so the
          exporter omits them.
    \item ``$r$ holds for $a$ and $b$'', a relation \emph{variable}, now counts as worded.
    \item The slow proof texts (\texttt{sup\_assoc}) got a time limit.
  \end{itemize}
  Result: 705 statements on four pages, 574 fully worded (\commit{976bbd0}).
\end{frame}

\begin{frame}{9. French}
  \prompt{Now, let us add French! Can you figure out what is needed?}
  \begin{columns}[T]
    \begin{column}{0.5\textwidth}
      Claude's findings:
      \begin{enumerate}
        \item the full grammar was six weeks old;
        \item French proof texts were \emph{empty}: \texttt{ProofUnitsFre}
              lacked the two functions that proof lines are made of;
        \item 27 words lacked French, in 5 modules;
        \item the page builder needed a French row;
        \item grammar flaws, for later.
      \end{enumerate}
      \prompt{Do 2-4 and add French as a third row for the time being}
    \end{column}
    \begin{column}{0.48\textwidth}
      \small
      \texttt{le\_antisymm}:\\
      Let $x$ be a set. Assume that $x$ is partially ordered. Let $a$ and $b$
      be elements of $x$. Assume that $a \leq b$. Assume that $b \leq a$.
      Then $a = b$.

      \vspace{1mm}
      Soit $x$ un ensemble. Supposons que $x$ est partiellement ordonné.
      Soient $a$ et $b$ des éléments de $x$. Supposons que $a \leq b$.
      Supposons que $b \leq a$. Alors $a = b$.
    \end{column}
  \end{columns}
\end{frame}

\begin{frame}{9. French: two interruptions}
  \small
  The first build took minutes: every French call loads the full grammar
  (2.8 s), and the proofs were read one at a time.
  \prompt{It is better to build an index that maps Dedukti to French and
  English and LaTeX in one go, so that the grammar is only loaded once. Then
  use this index to populate the html.}
  \begin{itemize}
    \item Informath (\commit{778e998}): \texttt{-languages=Eng,Fre} for
          \texttt{-json} and \texttt{-proof-text}, and \texttt{-no-back-to-dedukti},
          as the back-translation crashed on proofs.
    \item \texttt{build-index.py} $\to$ \texttt{index.json} in \textbf{27 s}, and
          \texttt{build-html.py} renders it in 0.06 s.
    \item Keyed by Dedukti names, the index revealed that the old pages had
          put some readings on the \bad{wrong declarations}.
  \end{itemize}
  \prompt{don't add german yet, just French and English}
  Claude had rebuilt the grammar with German, to keep a regression test
  running; the full grammar now has English and French only (\commit{823155a}).
\end{frame}

\begin{frame}{10. Czech}
  \prompt{Now that the grammar with 7 languages is in place, why not add Czech to the
  html file - to begin with. Please hide French and Czech by default but make them
  easy to click visible.}
  \begin{itemize}
    \item The Czech grammar had its user extensions switched off: \texttt{HottWordsCze}
          was a stub with two \bad{Swedish} words, and did not compile. Now in Czech,
          and on (\commit{0da10d2}).
    \item No \texttt{ProofUnitsCze}: proof texts would be empty, as French ones were.
    \item The same 27 words as for French: \texttt{MathlibCze} (\emph{nejmenší, shora
          omezený, množina dělitelů} \dots) and partial modules. Unreviewed.
    \item \texttt{build-index.py}: RunInformath loses its buffered output when it stops
          on an error, which cut off the proof \emph{before} a failing one; a proof now
          counts as failed only if it fails alone.
    \item A check box for each hidden language above the entries, remembered in the
          browser.
  \end{itemize}
  \prompt{add Czech and French links locally at each judgement, where they are easier
  to find.}
  A French and a Czech button in each entry's header line (\commit{0c9c416}).
\end{frame}

\begin{frame}{11. Aggregated hypotheses}
  \footnotesize
  \prompt{optimize hypotheses by using AdjKindHypo and SuchThatHypo \ldots\ even sorting
  the hypotheses so that this becomes more possible \ldots\ whenever the type
  dependencies are obeyed.}
  \begin{columns}[T]
    \begin{column}{0.5\textwidth}
      \texttt{aggregateHypos} in Informath (\commit{5907eec}):
      \begin{enumerate}
        \item one hypothesis per variable;
        \item a variable moves back next to the last one of its kind, if its kind
              depends on nothing in between;
        \item variables of a kind regrouped;
        \item the propositions about them attached: an adjective each, the same, or
              else ``such that''.
      \end{enumerate}
      The old rule had never applied: unnamed hypotheses arrive with generated labels.
    \end{column}
    \begin{column}{0.48\textwidth}
      \textbf{503 of 705} English readings change (\commit{afd0d9d}):
      \begin{itemize}
        \item 105 adjectives: ``Let $p$ and $q$ be \ok{prime natural numbers}.''
        \item 315 ``such that'': ``Let $p$ and $m$ be natural numbers \ok{such that}
              $p$ is prime and $m$ divides $p$'' --- $m$ was declared after ``Assume
              that $p$ is prime''.
        \item 83 class hypotheses: ``Let $x$ be a set such that $x$ is a join
              semilattice.''
      \end{itemize}
    \end{column}
  \end{columns}
  \textbf{Remaining:} \bad{``a set such that $x$ is a join semilattice''} could be ``a
  join semilattice''; ``such that'' after ``a function from \dots\ to sets'' seems to be
  about the sets; mixed wordings of a constant keep variables apart.
\end{frame}

\begin{frame}{12. German, and a time explosion}
  \small
  \prompt{I know that this can be heavy, but can you add German to Mathlib-informath?}
  \begin{itemize}
    \item German had the least: its user extensions were Naproche only, and its Wikidata
          lexicon partial. The table needed 67 German words, 40 of them general
          (\emph{Kardinalität, Bild, Vereinigung} \dots); new \texttt{ProofUnitsGer},
          \texttt{MathlibGer} and partial modules. Unreviewed.
    \item Then the batch call did not finish. GF's German linearization is exponential
          in nested noun phrases: ``the minimum of the minimum of $a$ and $b$ and $c$''
          --- depth 2 in 0.3 s, depth 3 in 35 s; English is instant.
  \end{itemize}
  \prompt{aggressive time limit per statement. I will try to fix German rgl.}
  \begin{itemize}
    \item German is read apart, in chunks of 12 with \textbf{1 s per statement}; a chunk
          over the limit is split in halves down to the slow statements, which get no
          German, marked on the page.
    \item 48 statements left out --- the lattice laws (\texttt{sup\_assoc} \dots) and
          the \texttt{ordCompl} lemmas; all 240 proofs have German. The index takes
          15 minutes instead of one.
    \item German flaws seen: \bad{``elemente''}, \bad{``Dann $2 \leq p$''} (without
          \emph{gilt}), ``eine prime natürliche Zahl'' for \emph{Primzahl}.
  \end{itemize}
\end{frame}

\begin{frame}[fragile]{13. Tactic proofs}
  \footnotesize
  \prompt{Now I'm interested in adding proofs that use tactics. This might be doable by
  ignoring proof expressions in Dedukti and going directly to lists of proof units.}
  \begin{columns}[T]
    \begin{column}{0.5\textwidth}
      242 of our 620 theorems are proved by tactics, mostly \texttt{rw}, \texttt{simp},
      \texttt{exact}; the scripts alone say little. Claude asked three questions:
      \ok{Lean's goals}, \ok{a new mode \texttt{-tactic-proof}}, \ok{Nat.Prime.Defs in
      English first}.
      \begin{itemize}
        \item \texttt{export --tactics}: the module elaborated again (5 s), each tactic
              of a sequence with its goals before and after, in Dedukti.
        \item \texttt{tactic-units.py}: units from comparing the goals --- assumed,
              obtained, stated, closed, ``it is enough to show'', cases.
        \item \texttt{RunInformath -tactic-proof}: the units through the symbol tables and
              the NLG, written out as by \texttt{-proof-text}.
      \end{itemize}
    \end{column}
    \begin{column}{0.48\textwidth}
\begin{lstlisting}
obtain ⟨n, hn⟩ := hm
have := pp.isUnit_or_isUnit hn
rw [Nat.isUnit_iff, Nat.isUnit_iff] at this
apply Or.imp_right _ this
rintro rfl
rw [hn, mul_one]
\end{lstlisting}
      \scriptsize
      \begin{tabular}{@{}r>{\raggedright\arraybackslash}p{0.88\linewidth}@{}}
        1. & Let $n$ be an element of $\mathbb{N}$ such that $p = m \times n$. \\
        2. & $m$ is a unit or $n$ is a unit. \\
        3. & $m = 1$ or $n = 1$. \\
        4. & It is enough to show that if $n = 1$, then $m = p$. \\
        5. & Assume that $n = 1$. \\
        6. & $m = p$. \\
      \end{tabular}
    \end{column}
  \end{columns}
  Fixed on the way: \texttt{rintro rfl} assumed the \bad{substituted} hypotheses;
  \texttt{by\_cases} lost its case hypotheses. The rebuild found that the RGL, changed in my
  German work, no longer compiles \bad{Czech}.
\end{frame}

\begin{frame}{14. Czech again, and optimized grammars}
  \small
  \prompt{can you fix the Czech grammar ?}
  \begin{itemize}
    \item The RGL's Czech had changed: \texttt{NP} with clitics and pro-drop, 32 agreement
          values, new \texttt{mkV2}, \texttt{mkVS}. Informath's \texttt{[ArgKind]} then
          had 536,870,912 concrete functions; an \texttt{ArgKindNP} without the new fields
          brings it to 1,048,576. ``Předpokládejme'' and ``všechna'' came for free.
    \item Then my \texttt{-optimize-pgf} for smaller grammars broke all proof texts:
          \bad{Index (-1) out of range}. Claude bisected it to the optimization; I knew
          why: units are not reachable from the start category, so I added \texttt{UnitJmt}.
    \item Claude added \texttt{ProofUnits} for Swedish, Czech and German (the units of
          tactic proofs), with ``dass $p$ prim \ok{ist}'' (verb last), and the Swedish
          words of the Mathlib pages: Swedish became the fifth language.
  \end{itemize}
\end{frame}

\begin{frame}{15. A symbol table for all of Mathlib}
  \small
  \prompt{Build a symbol table for all of Mathlib's definitions. Use information from
  docstrings and also try to parse camel script to multiword expressions. Add words to
  Informath lexicon if needed. Start by doing this just for English}
  \begin{itemize}
    \item \texttt{export --definitions}: all of Mathlib imported once (20 s), 56,867
          definitions, structures, classes, fields; 38,624 after leaving out tactics,
          notations, auxiliaries.
    \item \texttt{definition-table.py}: candidates from the docstring (``The cofinality of
          an ordinal is \dots'') and the name (\texttt{minFac}: ``the minimum factor of
          \#1''), Mathlib's abbreviations spelled out; all parsed in one call (30 s); the
          first whose category and arity fit is the entry.
    \item New words for the grammar from DictEng (categories, plurals): about 500.
    \item \textbf{70\%} of the definitions have a wording (26,154): functions 22,641 of
          31,858, types 1,792 of 2,457, predicates 1,721 of 4,309.
  \end{itemize}
\end{frame}

\begin{frame}{16. Readings that are right, or none}
  \footnotesize
  Rounds of sampling found the classes of errors:
  \begin{itemize}
    \item \texttt{AddMonoidHom}: \bad{``sum monoid homomorphism''} $\to$ \ok{``additive monoid
          homomorphism''}; \texttt{MulLeftMono} is about monotonicity, not monomorphisms.
    \item \texttt{Bicone.ofLimitCone}: \bad{``the limit cone of \#1''} $\to$ \ok{``the bicone
          of \#1''}; \texttt{pointOfClosedPoint} keeps its head, \texttt{PartitionOfUnity}
          its ``of''; \texttt{compRight} is the \ok{right composition}.
    \item Extra explicit arguments (a category, a topology): readings name the main ones,
          as \texttt{Eqt}'s ``\#2 is equal to \#3'' does.
    \item But \texttt{IsRelUpperSet s P} became \bad{``$P$ is a relation upper set''}. A
          wrong reading is worse than a symbolic one: a predicate of two objects is now
          binary, ``\#1 is a relative upper set with respect to \#2'', its subject the
          argument the docstring does not name after ``relative to'', ``for''; without such
          a clue, no reading (predicates 2,531 $\to$ 1,721).
  \end{itemize}
  The pages gain 14 fully worded statements: \texttt{UpperSet.carrier} is ``the carrier of
  $\mathrm{self}$''.
\end{frame}

\begin{frame}{17. The words in the other languages}
  \small
  \prompt{As for lexicon, you can also try ../wordnet. The abstract syntax file gives hints
  about what domain each word sense is used on.}
  \begin{itemize}
    \item GF's WordNet: 117 of 789 words, 13 by a sense tagged \emph{mathematics}, and
          even those mixed: \emph{Sechseck}, \emph{variance} \ok{right}; minus as
          \bad{\emph{Weg}}, triplet as \bad{\emph{trois}}; substructure as
          \bad{\emph{Infrastruktur}} by its gloss. Mathlib's jargon (cocone, hypercover,
          sheafification) is not there.
    \item Instead \texttt{translations.tsv}: one line per word with French, German,
          Swedish and Czech forms (gender, plural), written by Claude, \textbf{to be
          reviewed}; an empty cell where unsure. 473, 465, 441, 441 of 485 words.
    \item Where a language still lacks a word, \texttt{build-index.py} keeps the reading
          made without the generated table: 26 declarations, all with readings again.
    \item On the way: 300 junk words of docstrings out of the lexicon (\emph{because},
          \emph{banana}, \emph{simproc}).
  \end{itemize}
\end{frame}

\begin{frame}[fragile]{18. A stand-alone informalizer}
  \footnotesize
  \prompt{Write a stand-alone program that informalizes any piece of Lean code read from
  standard-in or from a file. \dots\ Claude would no more be needed when running this
  program.} \hfill (\texttt{standalone.md})
  \begin{columns}[T]
    \begin{column}{0.48\textwidth}
      \begin{itemize}
        \item \texttt{export --file}: any Lean file elaborated (about 35 s, mostly loading
              Mathlib), its declarations and tactic steps exported.
        \item \texttt{mathlib-informath}: the index code of the pages for one file;
              plain text, \texttt{-html}, \texttt{-pdf}; \texttt{-languages},
              \texttt{-table} for one's own entries; \texttt{-check}, \texttt{-install}.
        \item Later: a web server keeping Lean loaded, and VS Code.
      \end{itemize}
    \end{column}
    \begin{column}{0.5\textwidth}
\begin{lstlisting}
$ mathlib-informath Example.lean
prime_ne_one (theorem)
  English: Let $p$ be a prime natural number.
           Then $p \neq 1$.
  Proof:
    1. Assume that $p = 1$.     [intro h]
    2. $1$ is prime.            [rw [h] at hp]
    3. We have a contradiction.
                   [exact Nat.not_prime_one hp]
\end{lstlisting}
    \end{column}
  \end{columns}
\end{frame}

\begin{frame}{19. All definitions on a page}
  \small
  \prompt{I would like to have a browsable page with the definitions.}
  \begin{itemize}
    \item \texttt{definitions.html}: all 38,690 definitions with their wordings, 26,215 of
          them worded; the arguments by their names (``the cofinality of $o$''); the
          docstring's first sentence; a link to Mathlib's documentation.
    \item Search by name, wording and docstring; filters by part of Mathlib, by kind, and
          by where the wording comes from (by hand, docstring, name). The data is a 7 MB
          JSON file read by the page.
    \item Linked from the main page; the Dedukti rows of the entries are now folded, as
          the proofs are, as I asked: ``It is only interesting for debugging.''
  \end{itemize}
\end{frame}

\begin{frame}{20. All definitions read in English}
  \small
  \prompt{I would like to have one more thing in definitions.html: the informalizations of each
  definition. The docstrings are not needed}
  \begin{itemize}
    \item \texttt{export --definitions} also writes the definientia (25,494); all 38,690
          definitions read by the pages' code, 8 RunInformath processes, 1 s per definition:
          \textbf{half an hour}.
    \item \textbf{38,279} (99\%) with a reading, 405 over the time limit; about 12,300 (32\%)
          with no constant shown symbolically.
    \item A sample of 99 first, as I asked; but it loaded its data by \bad{\texttt{fetch}}, which
          a browser refuses for a file: \prompt{I couldn't see anything. Why not build a static html
          similar to mathlib-informath?}
    \item Now 111 static pages of at most 400 definitions, and a search that jumps to them.
    \item ``Let $V$ be a set. Then $V$ is thin, if for all elements $a$ and $b$ of $V$, the set
          of homomorphisms of $a$ and $b$ is a subsingleton.'' (\texttt{Quiver.IsThin})
    \item The other languages: possible on demand, same pipeline; not made yet.
  \end{itemize}
\end{frame}

\begin{frame}{The result}
  \centering
  Fully worded statements:

  \begin{tabular}{@{}lrrrrrr@{}}
    \toprule
    page & statements & English & French & Czech & German & Swedish \\
    \midrule
    Mathlib.Order.Defs & 354 & 291 & 289 & 289 & 248 & 289 \\
    Mathlib.Algebra & 29 & 29 & 29 & 29 & 28 & 28 \\
    Mathlib.Data.Nat.Prime & 128 & 125 & 125 & 125 & 125 & 125 \\
    Mathlib.Data.Nat.Factorization & 194 & 143 & 143 & 143 & 132 & 142 \\
    \midrule
    all & \textbf{705} & \textbf{588} & \textbf{586} & \textbf{586} & \textbf{533} & \textbf{584} \\
    \bottomrule
  \end{tabular}

  \vspace{3mm}
  \raggedright
  \begin{itemize}
    \item 264 declarations with their definiens; 240 proof texts in each
          language, out of 242 term proofs; 51 statements without German (time limit).
    \item 25 tactic proofs in all five languages (\texttt{Nat.Prime.Defs}), out of 242.
    \item 25,857 of Mathlib's 38,690 definitions worded; 38,279 read in English
          (\texttt{definitions.html}).
    \item Every statement has a formula; 31 need \texttt{-force-symbolic}.
    \item \url{https://aarneranta.github.io/mathlib-informath/mathlib-informath.html}
  \end{itemize}
\end{frame}

% ------------------------------------------------------------------ open problems

\section{Open problems}

\begin{frame}{Open problems: coverage}
  \small
  \begin{itemize}
    \item \textbf{Scale.} 705 of Mathlib's 240,000 declarations on the pages; 26,154
          generated entries for its definitions, unreviewed: composite names
          (\texttt{XEquivY}), bundled maps without arguments (\bad{``the dual''}), verbs
          (\texttt{PreservesX}), 2,000 unknown abbreviations; the words of the other
          languages in \texttt{translations.tsv} to be reviewed.
    \item \textbf{Tactic proofs.} 242 of the 620 theorems here, and most of
          Mathlib's, are proved \texttt{by} tactics: 25 have proof texts so far, in
          English; the others, the other languages, and the term and tactic proofs to be
          told apart on the page. Long \texttt{simp\_rw} goals, \texttt{by\_contra} as
          ``it is enough to show that we have a contradiction''.
    \item \textbf{Unworded statements} (131 of 705):
      \begin{itemize}
        \item multi-letter variables (\texttt{lt}, \texttt{self},
              \texttt{motive}) are shown as $\mathrm{lt}$ and counted as constants;
        \item function arguments: \texttt{Set.ofPred}, the motives of induction;
        \item \texttt{Leq} used unapplied: $\mathrm{Leq}(x_2, x_3)$.
      \end{itemize}
    \item Decidability instances defined as declarations are still exported;
          two \texttt{Decidable} proofs (with \texttt{dite}) have no proof text.
    \item Theorem labels are mangled names: \texttt{Nat\_\_Prime\_\_two\_le}.
  \end{itemize}
\end{frame}

\begin{frame}{Open problems: language}
  \scriptsize
  \begin{itemize}
    \item \textbf{Proof texts} where a lemma or constructor is used unapplied:
          ``For all sets $x$ and functions $r$ \ldots\ and elements
          $\mathrm{irrefl}$ of $\mathrm{Proof}$ applied to the proposition
          that \ldots, $r$ is irreflexive.'' Bound variables are renumbered
          ($x_4$, $x_5$) where the source reuses names.
    \item \textbf{Meaning that depends on the type}: $\leq$ on finite sets reads
          ``less than or equal to'', where Mathlib means ``subset of''.
    \item \textbf{French}: ``pour tous \bad{des} nombres naturels'', ``il
          n'est pas le cas que'' (a calque of ``it is not the case that''),
          ``une relation \bad{d'équivalences}''; three French modules are partial.
    \item \textbf{Czech}: unreviewed words; ``existuje přirozené číslo $p$,
          \bad{takový}, že'' (\emph{takové}): ``such that'' agrees only after a CN.
    \item \textbf{Informath}:
      \begin{itemize}
        \item \texttt{\$} in a quoted verbal form splits a table entry
              (\texttt{splitEntry}), which rules out binder shapes in tables;
        \item a proposition inside a macro argument loses its notation;
        \item ``does not have \bad{a} upper bound'';
        \item \texttt{proof2dedukti} does not cover proofs, so the
              back-translation of \texttt{-json} crashes on them.
      \end{itemize}
    \item \textbf{German}: GF's linearization is exponential in nested noun phrases
          (48 statements over the limit); ``elemente''; ``Dann $2 \leq p$''.
    \item \textbf{Swedish}: ``lika \bad{till}'' (\emph{lika med}) and Czech ``rovná
          \bad{do}'': \texttt{Eqt}'s ``is equal to'' parses word by word.
  \end{itemize}
\end{frame}

% ------------------------------------------------------------------ summary

\section{Summary}

\begin{frame}{Who did what}
  \scriptsize
  \renewcommand{\arraystretch}{1.0}\begin{tabularx}{\textwidth}{@{}>{\raggedright\arraybackslash}X>{\raggedright\arraybackslash}X>{\raggedright\arraybackslash}X@{}}
    \toprule
    \textbf{Already in Informath} & \textbf{From my prompts} & \textbf{Added by Claude} \\
    \midrule
    The GF grammar and its RGL-based concrete syntaxes, French among them
    (2025--; \texttt{WikidataWordsFre} completed with Claude in July)
    &
    The goal: an informalization system for Mathlib with Informath
    &
    \texttt{Export.lean}: the translation rules, classes as predicates,
    \texttt{\_\_app} coercions, eta-expanded binders, proofs in \texttt{nd.dk}
    rules, filters \\
    The Dedukti dialect, its base constants, symbol tables with verbal,
    symbolic and raw macro forms
    &
    The choices among Claude's options: own exporter, Nat primes first, a new
    project with the cache
    &
    \texttt{mapping.tsv}; \texttt{mathlib.dkgf} (108 entries) with
    \texttt{suggest-entries.py} \\
    NLG with variants and ranking; \texttt{-to-symbolic-latex},
    \texttt{-force-symbolic}, \texttt{-parse-example}
    &
    The page's contents; GitHub hosting (I created the repository and pushed)
    &
    \texttt{build-index.py}, \texttt{build-html.py}: index, pages, macros,
    proof links \\
    \texttt{-proof-text} (since June, 8 of 21 commits with Claude), and
    godement-index's way of showing proofs
    &
    A page per part; keeping the definientia and showing the proofs;
    words from the sources; new words in \texttt{Mathlib.gf}
    &
    Words in \texttt{Mathlib.gf}, \texttt{MathlibFre.gf}, and partial French
    modules \\
    The regression suite, the autoformalization template
    &
    The next modules; French as a third row; one index loading the grammar
    once; Czech, hidden; hypotheses aggregated, reordered; German, with an aggressive time limit;
    tactic proofs as proof units, after Verbose Lean
    &
    Informath: \texttt{abstractionTypes}, \texttt{-languages},
    multilingual proof texts, \texttt{-no-back-to-dedukti},
    \texttt{ProofUnitsFre}, Czech, \texttt{aggregateHypos}, \texttt{-tactic-proof}; \texttt{export --tactics} \\
    \texttt{-optimize-pgf}, \texttt{UnitJmt}; the Czech RGL
    &
    Swedish; a symbol table of all definitions; WordNet; a reviewable table; older
    readings kept; the stand-alone program (\texttt{standalone.md}); the definitions page
    &
    \texttt{export --definitions, --file}; \texttt{definition-table.py}, 500 words and
    their translations; \texttt{mathlib-informath}; \texttt{definitions.html} \\
    \bottomrule
  \end{tabularx}

  \vspace{1mm}
  Claude also found and reported the bugs along the way, its own among them
  (the red macros, the misaligned readings), and asked before committing; I pushed.
\end{frame}

\begin{frame}[standout]
  \url{github.com/aarneranta/mathlib-informath}

  \vspace{3mm}
  \normalsize
  705 Mathlib statements, 265 proofs, in English, French, Czech, German, Swedish and LaTeX;
  all of Mathlib's 38,690 definitions read in English
\end{frame}

\end{document}
